% Part: second-order-logic
% Chapter: syntax-and-semantics
% Section: satisfaction

\documentclass[../../../include/open-logic-section]{subfiles}

\begin{document}

\olfileid{sol}{syn}{sat}

\olsection{पूर्ति}

\begin{explain}
द्वितीय-क्रम
!!{formula}s साठी
पूर्ति-संबंध
$\Sat{M}{!A}[s]$
परिभाषित करताना
द्वितीय-क्रम
!!{variable}s
समाविष्ट करण्यासाठी
व्याख्या वाढवाव्या
लागतात.
!!a{structure}
ही संकल्पना
द्वितीय-क्रम
आणि प्रथम-क्रम
तर्कशास्त्रात
सारखीच असते.
फरक फक्त
चर-मूल्यांकन~$s$
बाबत असतो:
आता त्याने
प्रथम-क्रम
!!{variable}s
बरोबरच
द्वितीय-क्रम
!!{variable}s
यांनाही मूल्ये
द्यावी लागतात.
\end{explain}

\begin{defn}[चर-मूल्यांकन]
!!a{structure}~$\Struct{M}$
साठीचे
\emph{चर-मूल्यांकन}~$s$
हे असे फलन
आहे जे पुढीलपैकी
प्रत्येकाला
दिल्याप्रमाणे
प्रतिचित्रित करते:
\begin{enumerate}
\item वस्तू-!!{variable}~$\Obj{v_i}$
  याला~$\Domain M$
  मधील घटकाकडे,
  म्हणजे
  $s(\Obj{v_i}) \in \Domain{M}$;
\item $n$-स्थानी
  संबंध-चल~$\Obj{V_i^n}$
  याला~$\Domain{M}$
  वरील $n$-स्थानी
  संबंधाकडे,
  म्हणजे
  $s(\Obj{V_i^n}) \subseteq \Domain{M}^n$;
\item $n$-स्थानी
  फलन-चल~$\Obj{u_i^n}$
  याला
  $\Domain{M}$
  वरून $\Domain{M}$
  मध्ये जाणाऱ्या
  $n$-स्थानी
  फलनाकडे,
  म्हणजे
  $s(\Obj{u_i^n})\colon \Domain{M}^n \to \Domain{M}$;
\end{enumerate}
\end{defn}

\begin{explain}
!!^a{structure}
प्रत्येक !!{constant}
आणि !!{function}
ला !!a{value}
देते, तर द्वितीय-क्रम
चर-मूल्यांकन
प्रत्येक वस्तू-चल
आणि फलन-चल
यांना अनुक्रमे
वस्तू आणि फलने
देते. या दोन्हींच्या
साहाय्याने
प्रत्येक पदाला
मूल्य देता येते.
\end{explain}

\begin{defn}[पदाचे \usetoken{S}{value}]
$t$ हे भाषा~$\Lang L$
मधील पद,
$\Struct M$ ही~$\Lang L$
साठी !!a{structure},
आणि $s$ हे~$\Struct M$
साठी !!a{variable}
मूल्यांकन असेल,
तर \emph{!!{value}}~
$\Value{t}{M}[s]$
हे प्रथम-क्रम
पदांप्रमाणेच,
पुढील अतिरिक्त
खंडासह परिभाषित
होते:
\begin{quote}
\indcase{t}{\Atom{u}{t_1, \ldots, t_n}}{
\[
\Value{\indfrm}{M}[s] = s(u)(\Value{t_1}{M}[s], \ldots,
\Value{t_n}{M}[s]).
\]}
\end{quote}
\end{defn}

\begin{defn}[$x$-पर्याय]
$s$ हे !!a{structure}~$\Struct M$
साठी !!a{variable}
मूल्यांकन असेल,
तर~$\Struct M$
साठीचे असे
कोणतेही !!{variable}
मूल्यांकन $s'$
जे $s$ पेक्षा
जास्तीत जास्त
$x$ ला दिलेल्या
मूल्यातच भिन्न
असेल, त्याला
$s$ चा
\emph{$x$-पर्याय}
म्हणतात.
$s'$ हा $s$
चा $x$-पर्याय
असेल, तर
$\varAssign{s'}{s}{x}$
असे लिहितो.
(द्वितीय-क्रम
चल $X$ किंवा
$u$ यांसाठीही
तसेच.)
\end{defn}

\begin{defn}
  %READERNOTE{SOL6-A091: पर्याय मूल्यांकन निश्चित x, X किंवा u वर बदलते; त्याचे कोणत्याही y वरील मूल्य डाव्या बाजूला दाखवले आहे. पूरक संचासाठी disjointness बरोबर पूर्ण आधारक्षेत्र झाकणे आवश्यक आहे. Conjunction encoding मधील X मुक्त नसण्याची अट capture टाळते.}
  $s$ हे !!a{structure}~$\Struct M$
  साठी !!a{variable}
  मूल्यांकन
  आणि $m \in \Domain{M}$
  असेल, तर
  मूल्यांकन~$\Subst{s}{m}{x}$
  हे पुढीलप्रमाणे
  परिभाषित केलेले
  !!{variable}
  मूल्यांकन आहे:
  \[\Subst{s}{m}{x}(y) = \begin{cases}
    m & \text{जर } y \ident x\\
    s(y) & \text{अन्यथा},
  \end{cases}\]
  $X$ हे $n$-स्थानी
  संबंध-!!{variable}
  आणि $M \subseteq
  \Domain{M}^n$
  असेल, तर
  $\Subst{s}{M}{X}$
  हे पुढीलप्रमाणे
  परिभाषित केलेले
  !!{variable}
  मूल्यांकन आहे:
  \[\Subst{s}{M}{X}(y) = \begin{cases}
    M & \text{जर } y \ident X\\
    s(y) & \text{अन्यथा}.
  \end{cases}\]
  $u$ हे $n$-स्थानी
  फलन-!!{variable}
  आणि $f\colon \Domain{M}^n
  \to \Domain{M}$
  असेल, तर
  $\Subst{s}{f}{u}$
  हे पुढीलप्रमाणे
  परिभाषित केलेले
  !!{variable}
  मूल्यांकन आहे:
  \[\Subst{s}{f}{u}(y) = \begin{cases}
    f & \text{जर } y \ident u\\
    s(y) & \text{अन्यथा}.
  \end{cases}\]
  प्रत्येक बाबतीत
  $y$ हे कोणतेही
  प्रथम-क्रम किंवा
  द्वितीय-क्रम
  !!{variable}
  असू शकते.
\end{defn}

\begin{defn}[पूर्ति]
द्वितीय-क्रम
!!{formula}s~$!A$
साठी पूर्तीची
व्याख्या
\olref[fol][syn][sat]{defn:satisfaction}
प्रमाणेच आहे;
त्यात पुढील
खंड वाढवले
आहेत:
\begin{enumerate}
\item \indcase{!A}{\Atom{X^n}{t_1, \dots, t_n}}{$\Sat{M}{\indfrm}[s]$
  हे $\langle \Value{t_1}{M}[s], \dots, \Value{t_n}{M}[s] \rangle \in
  s(X^n)$ असेल तर आणि तरच खरे.}
\tagitem{prvAll}{%
  \indcase{!A}{\lforall[X][!B]}{$\Sat{M}{\indfrm}[s]$ हे प्रत्येक $M
  \subseteq \Domain{M}^n$ साठी $\Sat{M}{!B}[\Subst{s}{M}{X}]$ असेल
  तर आणि तरच खरे.}}{}

\tagitem{prvEx}{%
  \indcase{!A}{\lexists[X][!B]}{$\Sat{M}{\indfrm}[s]$ हे किमान
  एका $M \subseteq \Domain{M}^n$ साठी
  $\Sat{M}{!B}[\Subst{s}{M}{X}]$ असेल तर आणि तरच खरे.}}{}

\tagitem{prvAll}{%
  \indcase{!A}{\lforall[u][!B]}{$\Sat{M}{\indfrm}[s]$ हे प्रत्येक
  $f\colon \Domain{M}^n \to \Domain{M}$ साठी
  $\Sat{M}{!B}[\Subst{s}{f}{u}]$ असेल तर आणि तरच खरे.}}{}

\tagitem{prvEx}{%
  \indcase{!A}{\lexists[u][!B]}{$\Sat{M}{\indfrm}[s]$ हे किमान
  एका $f\colon \Domain{M}^n \to \Domain{M}$ साठी
  $\Sat{M}{!B}[\Subst{s}{f}{u}]$ असेल तर आणि तरच खरे.}}{}
\end{enumerate}
\end{defn}

\begin{ex}
  !!{formula}
  $\lforall[z][(\Atom{X}{z} \liff \lnot
    \Atom{Y}{z})]$
  विचारात घ्या.
  त्यात द्वितीय-क्रम
  संख्यक नाहीत,
  पण द्वितीय-क्रम
  !!{variable}s
  $X$ आणि~$Y$
  आहेत (येथे ती
  एक-स्थानी मानली
  आहेत).
  याच्याशी संबंधित
  प्रथम-क्रम
  !!{sentence}
  $\lforall[z][(\Atom{P}{z} \liff \lnot \Atom{R}{z})]$
  असे म्हणते की
  आधारक्षेत्रातील प्रत्येक वस्तू
  $P$ आणि $R$ यांच्या
  अर्थनिर्धारणांपैकी नेमक्या एकात येते:
  ती पहिल्यात येते, तर आणि तरच
  ती दुसऱ्यात येत नाही.
  !!a{structure}
  मध्ये
  !!a{predicate}~$P$
  चे अर्थनिर्धारण
  $\Assign{P}{M}$
  याने दिलेले
  असते. पण
  $X$ आणि~$Y$
  सारख्या द्वितीय-क्रम
  !!{variable}s
  चे अर्थनिर्धारण
  !!{structure}
  स्वतः देत नाही;
  ते !!a{variable}
  मूल्यांकन देते.
  हे द्वितीय-क्रम
  !!{formula}
  !!a{sentence}
  नाही (त्यात
  $X$ आणि~$Y$
  ही मुक्त
  !!{variable}s
  आहेत);
  म्हणून त्याची
  पूर्ति फक्त
  !!a{structure}~$\Struct{M}$
  आणि त्यासोबत
  !!a{variable}
  मूल्यांकन~$s$
  यांच्या सापेक्ष
  ठरते.

  $\Sat{M}{\lforall[z][(\Atom{X}{z} \liff \lnot \Atom{Y}{z})]}[s]$
  हे $s(X)$ मधील
  !!{element}s आणि $s(Y)$ मधील
  !!{element}s एकमेकांपासून वेगळे असतील
  आणि दोन्ही संच मिळून संपूर्ण आधारक्षेत्र झाकतील,
  तर आणि तरच खरे असते; म्हणजे
  $s(Y) = \Domain{M} \setminus s(X)$
  असेल तर आणि
  तरच खरे असते.
  उदाहरणार्थ,
  $\Domain{M} = \{1, 2, 3\}$
  घ्या.
  येथे !!{predicate}s,
  !!{function}s
  किंवा !!{constant}s
  येत नसल्यामुळे
  $\Struct{M}$
  चे आधारक्षेत्र
  एवढेच संबंधित
  आहे.
  आता
  $s_1(X) = \{1, 2\}$
  आणि
  $s_1(Y) = \{3\}$
  असताना
  $\Sat{M}{\lforall[z][(\Atom{X}{z}
      \liff \lnot \Atom{Y}{z})]}[s_1]$
  खरे आहे.

  याउलट,
  $s_2(X) = \{1, 2\}$
  आणि
  $s_2(Y) = \{2, 3\}$
  असल्यास
  $\Sat/{M}{\lforall[z][(\Atom{X}{z} \liff \lnot
      \Atom{Y}{z})]}[s_2]$.
  कारण
  $\Sat{M}{\Atom{X}{z}}[\Subst{s_2}{2}{z}]$
  (कारण
  $2 \in \Subst{s_2}{2}{z}(X)$),
  पण
  $\Sat/{M}{\lnot \Atom{Y}{z}}[\Subst{s_2}{2}{z}]$
  (कारण
  $2 \in
  \Subst{s_2}{2}{z}(Y)$
  सुद्धा).
\end{ex}

\begin{ex}
  $\Sat{M}{\lexists[Y][(\lexists[y][\Atom{Y}{y}] \land
  \lforall[z][(\Atom{X}{z} \liff \lnot \Atom{Y}{z})])]}[s]$
  हे तेव्हाच खरे
  असते जेव्हा
  असा $N \subseteq \Domain{M}$
  असतो की
  $\Sat{M}{(\lexists[y][\Atom{Y}{y}] \land \lforall[z][(\Atom{X}{z}
  \liff \lnot \Atom{Y}{z})])}[\Subst{s}{N}{Y}]$.
  आणि असे
  $N \neq \emptyset$
  असताना
  ($\Sat{M}{\lexists[y][\Atom{Y}{y}]}[\Subst{s}{N}{Y}]$
  मिळण्यासाठी)
  आणि आधीच्या
  उदाहरणाप्रमाणे
  $N = \Domain{M} \setminus s(X)$
  असताना घडते.
  दुसऱ्या शब्दांत,
  $\Sat{M}{\lexists[Y][(\lexists[y][\Atom{Y}{y}] \land
  \lforall[z][(\Atom{X}{z} \liff \lnot \Atom{Y}{z})])]}[s]$
  हे
  $\Domain{M} \setminus s(X)$
  रिकामे नसेल,
  म्हणजे
  $s(X) \neq
  \Domain{M}$
  असेल, तर आणि
  तरच खरे असते.
  म्हणून
  $\Domain{M}
  = \{1, 2, 3\}$
  आणि
  $s(X) = \{1, 2\}$
  असल्यास
  !!{formula}
  ची पूर्ति होते;
  पण
  $s(X) = \{1, 2, 3\}
  = \Domain{M}$
  असल्यास होत नाही.

  $s(X) = \Domain{M}$
  असताना या
  !!{formula}
  ची पूर्ति
  होत नसल्यामुळे,
  पुढील
  !!{sentence}
  \[
  \lforall[X][\lexists[Y][(\lexists[y][\Atom{Y}{y}] \land
      \lforall[z][(\Atom{X}{z} \liff \lnot \Atom{Y}{z})])]]
  \]
  कधीच पूर्ण
  होत नाही:
  कोणतीही
  !!{structure}~$\Struct{M}$
  घेतली तरी
  $s(X) = \Domain{M}$
  हे मूल्यांकन
  त्या
  !!{sentence}
  ला असत्य
  ठरवेल.
  याउलट,
  पुढील वाक्य
  \[
  \lexists[X][\lexists[Y][(\lexists[y][\Atom{Y}{y}] \land
      \lforall[z][(\Atom{X}{z} \liff \lnot \Atom{Y}{z})])]]
  \]
  कोणत्याही
  मूल्यांकन~$s$
  च्या सापेक्ष
  पूर्ण होते,
  कारण
  $M \subseteq \Domain{M}$
  पण
  $M \neq \Domain{M}$
  असा संच नेहमी
  निवडता येतो
  (उदा.,
  $M = \emptyset$).
\end{ex}

\begin{ex}
  द्वितीय-क्रम
  !!{sentence}~
  $\lforall[X][\lforall[y][X(y)]]$
  असे म्हणते की
  प्रत्येक
  $1$-स्थानी
  संबंध, म्हणजे
  प्रत्येक
  गुणधर्म, प्रत्येक
  वस्तूला लागू
  होतो.
  हे स्पष्टपणे
  कधीच खरे नसते:
  प्रत्येक~$\Struct{M}$
  मध्ये
  $s(X) = \emptyset$
  आणि
  $s(y) = a \in
  \Domain{M}$
  असे चर-मूल्यांकन~$s$
  घेतल्यास
  $\Sat/{M}{X(y)}[s]$
  मिळते.
  याचा अर्थ
  $!A \lif
  \lforall[X][\lforall[y][X(y)]]$
  हे द्वितीय-क्रम
  तर्कशास्त्रात
  $\lnot !A$
  शी सममूल्य
  आहे; म्हणजे
  $\Sat{M}{!A \lif
  \lforall[X][\lforall[y][X(y)]]}$
  हे
  $\Sat{M}{\lnot !A}$
  असेल तर आणि
  तरच खरे.
  दुसऱ्या शब्दांत,
  द्वितीय-क्रम
  तर्कशास्त्रात
  $\lforall$
  आणि~$\lif$
  वापरून
  $\lnot$
  परिभाषित
  करता येते.
\end{ex}

\begin{prob}
  द्वितीय-क्रम
  तर्कशास्त्रात
  $\lforall$
  आणि~$\lif$
  वापरून इतर
  संयोजके
  परिभाषित
  करता येतात
  हे दाखवा:
  \begin{enumerate}
    \item येथे X हे संयोजित करावयाच्या दोन्ही सूत्रांमध्ये मुक्त नसलेले
    एकस्थानी विधेय चल निवडा. द्वितीय-क्रम तर्कशास्त्रात $!A \land !B$ हे
    $\lforall[X][(!A \lif (!B \lif \lforall[x][X(x)])\lif
    \lforall[x][X(x)])]$ शी सममूल्य आहे हे सिद्ध करा.
    \item फक्त $\lforall$ आणि $\lif$ वापरून
    $!A \lor !B$ शी सममूल्य द्वितीय-क्रम सूत्र शोधा.
  \end{enumerate}
\end{prob}

\end{document}
